\chapter{Fundamental Concepts and Solution Chemistry}
\label{chap:exo_ch01}

\begin{quicksummary}{Core Definitions, Concentration Scales \& Dilution Laws}
\begin{enumerate}[leftmargin=1.4em,itemsep=3pt]
  \item \textbf{Amount of Substance (Mole) \& Avogadro's Number:}
  \[
  n = \frac{m}{M} = \frac{N}{N_A} = \frac{V_{\text{gas}}}{V_0} \quad \text{where } N_A = 6.022\,141 \times 10^{23}\ \mathrm{mol^{-1}},\ V_0 = 22.414\ \mathrm{L\cdot mol^{-1}}\ \text{(STP)}.
  \]
  \item \textbf{Molarity ($C_M$ or $[\ce{X}]$) \& Mass Concentration ($C_m$):}
  \[
  C_M = \frac{n_{\text{solute}}}{V_{\text{sol}}} \quad [\mathrm{mol\cdot L^{-1}} \equiv \mathrm{M}], \qquad C_m = \frac{m_{\text{solute}}}{V_{\text{sol}}} \quad [\mathrm{g\cdot L^{-1}}], \qquad C_m = C_M \cdot M.
  \]
  \item \textbf{Molality ($m$ or $b$) \& Commercial Reagent Relation:}
  \[
  m = \frac{n_{\text{solute}}}{m_{\text{solvent}}\ (\mathrm{kg})} \quad [\mathrm{mol\cdot kg^{-1}}], \qquad C_M = \frac{10 \cdot P\% \cdot d}{M} \quad (d = \rho_{\text{sol}}/\rho_{\text{water}}).
  \]
  \item \textbf{Normality ($N$) \& Equivalent Mass ($E_q$):}
  \[
  N = Z \cdot C_M = \frac{m_{\text{solute}}}{E_q \cdot V_{\text{sol}}}, \qquad E_q = \frac{M}{Z}.
  \]
  \begin{itemize}[leftmargin=1.2em,itemsep=1pt]
    \item \textit{Acids/Bases:} $Z = \text{number of exchangeable } \ce{H+} \text{ or } \ce{OH-}$.
    \item \textit{Salts/Precipitates:} $Z = \text{total positive charge of metal cations} = n_{\text{cation}} \times z^+$.
    \item \textit{Redox:} $Z = \text{number of electrons transferred per mole of reagent}$.
  \end{itemize}
  \item \textbf{Dilution Law \& Titration Equivalence Principle:}
  \[
  n_{\text{initial}} = n_{\text{final}} \implies C_1 V_1 = C_2 V_2, \qquad \text{Equivalence:} \quad N_1 V_1 = N_2 V_2.
  \]
  \item \textbf{Mixing Non-reacting Solutions of the Same Solute:}
  \[
  C_{\text{final}} = \frac{n_1 + n_2 + \dots}{V_1 + V_2 + \dots} = \frac{C_1 V_1 + C_2 V_2 + \dots}{V_{\text{total}}}.
  \]
\end{enumerate}
\end{quicksummary}

\vspace{10pt}

% ===========================================================================
% EXERCISE 1.1
% ===========================================================================
\begin{exercise}[$\bigstar$ Foundational]{States of Matter and Phase Transition Enthalpy}
Classify the following physical transformations according to their thermodynamic nature (endothermic or exothermic), state the phase transition name, and describe the molecular-level change in kinetic and potential energy:
\begin{enumerate}[leftmargin=1.5em,itemsep=2pt]
  \item Formation of frost on a cold winter morning directly from atmospheric water vapor.
  \item The disappearance of dry ice ($\ce{CO2(s)}$) placed on a laboratory bench at room temperature without liquid formation.
  \item The condensation of steam on a chilled condenser coil during a distillation experiment.
\end{enumerate}
\end{exercise}

\begin{tabletworkspace}[5.0cm]
\end{tabletworkspace}

\begin{solution}
\textbf{1. Frost Formation (Water vapor $\longrightarrow$ Solid ice):}
\begin{itemize}[leftmargin=1.5em,itemsep=2pt]
  \item \textbf{Transition Name:} \textbf{Deposition} (or desublimation / solid condensation).
  \item \textbf{Thermodynamic Nature:} \textbf{Exothermic} ($\Delta H_{\text{dep}} < 0$). Heat is released into the surroundings.
  \item \textbf{Microscopic Description:} High-energy, disordered gaseous water molecules lose kinetic energy upon contact with the cold surface. Intermolecular hydrogen bonds form, locking water molecules into an ordered crystalline lattice with minimal potential energy.
\end{itemize}

\textbf{2. Disappearance of Dry Ice ($\ce{CO2(s)} \longrightarrow \ce{CO2(g)}$):}
\begin{itemize}[leftmargin=1.5em,itemsep=2pt]
  \item \textbf{Transition Name:} \textbf{Sublimation}.
  \item \textbf{Thermodynamic Nature:} \textbf{Endothermic} ($\Delta H_{\text{sub}} > 0$). Heat is absorbed from the room.
  \item \textbf{Microscopic Description:} Thermal energy from the room overcomes the weak van der Waals dispersion forces holding non-polar $\ce{CO2}$ molecules in the solid lattice, increasing both intermolecular separation (potential energy) and molecular velocity (kinetic energy).
\end{itemize}

\textbf{3. Condensation of Steam ($\ce{H2O(g)} \longrightarrow \ce{H2O(l)}$):}
\begin{itemize}[leftmargin=1.5em,itemsep=2pt]
  \item \textbf{Transition Name:} \textbf{Condensation} (or liquefaction).
  \item \textbf{Thermodynamic Nature:} \textbf{Exothermic} ($\Delta H_{\text{cond}} = -\Delta H_{\text{vap}} < 0$).
  \item \textbf{Microscopic Description:} Heat is transferred to the condenser coolant, decreasing molecular velocities. Attractive dipole-dipole and hydrogen-bonding interactions pull vapor molecules together into a dense, fluid condensed liquid phase.
\end{itemize}
\end{solution}

% ===========================================================================
% EXERCISE 1.2
% ===========================================================================
\begin{exercise}[$\bigstar$ Foundational]{The Mole Concept and Avogadro Counting}
A pure sample of metallic iron ($\ce{Fe}$) has a mass of $m = 16.755\ \mathrm{g}$.
\begin{enumerate}[leftmargin=1.5em,itemsep=2pt]
  \item Calculate the amount of substance $n$ of iron in the sample ($M(\ce{Fe}) = 55.845\ \mathrm{g\cdot mol^{-1}}$).
  \item Compute the absolute number of iron atoms $N_{\text{atoms}}$ contained in this sample using $N_A = 6.022\,141 \times 10^{23}\ \mathrm{mol^{-1}}$.
  \item What mass of elemental sulfur ($\ce{S}$, $M = 32.065\ \mathrm{g\cdot mol^{-1}}$) contains exactly the same number of atoms as this iron sample?
\end{enumerate}
\end{exercise}

\begin{tabletworkspace}[5.0cm]
\end{tabletworkspace}

\begin{solution}
\textbf{1. Amount of Substance of Iron ($n_{\text{Fe}}$):}
\[
n_{\text{Fe}} = \frac{m_{\text{Fe}}}{M(\ce{Fe})} = \frac{16.755\ \mathrm{g}}{55.845\ \mathrm{g\cdot mol^{-1}}} = \mathbf{0.3000\ \mathrm{mol}}.
\]

\textbf{2. Absolute Number of Iron Atoms ($N_{\text{atoms}}$):}
\[
N_{\text{atoms}} = n_{\text{Fe}} \cdot N_A = 0.3000\ \mathrm{mol} \times 6.022\,141 \times 10^{23}\ \mathrm{mol^{-1}} = \mathbf{1.8066 \times 10^{23}\ \text{atoms}}.
\]

\textbf{3. Equivalent Mass of Sulfur:}
To contain identical numbers of atoms, the amounts of substance must be identical: $n_{\text{S}} = n_{\text{Fe}} = 0.3000\ \mathrm{mol}$.
\[
m_{\text{S}} = n_{\text{S}} \cdot M(\ce{S}) = 0.3000\ \mathrm{mol} \times 32.065\ \mathrm{g\cdot mol^{-1}} = \mathbf{9.6195\ \mathrm{g}}.
\]
\end{solution}

% ===========================================================================
% EXERCISE 1.3
% ===========================================================================
\begin{exercise}[$\bigstar$ Foundational]{Molar Mass of Hydrated and Polyatomic Crystals}
Calculate the molar mass of the following chemical compounds to two decimal places:
\begin{enumerate}[leftmargin=1.5em,itemsep=2pt]
  \item Anhydrous calcium phosphate: $\ce{Ca3(PO4)2}$.
  \item Copper(II) sulfate pentahydrate: $\ce{CuSO4*5H2O}$.
  \item Potassium permanganate: $\ce{KMnO4}$.
  \item Ammonium iron(II) sulfate hexahydrate (Mohr's Salt): $(\ce{NH4})_2\ce{Fe(SO4)2*6H2O}$.
\end{enumerate}
\textit{Standard atomic masses ($\mathrm{g\cdot mol^{-1}}$):} $\ce{H} = 1.008$, $\ce{N} = 14.007$, $\ce{O} = 15.999$, $\ce{P} = 30.974$, $\ce{S} = 32.065$, $\ce{K} = 39.098$, $\ce{Ca} = 40.078$, $\ce{Mn} = 54.938$, $\ce{Fe} = 55.845$, $\ce{Cu} = 63.546$.
\end{exercise}

\begin{tabletworkspace}[5.0cm]
\end{tabletworkspace}

\begin{solution}
\textbf{1. Calcium phosphate $\ce{Ca3(PO4)2}$:}
\begin{align*}
M &= 3 \cdot M(\ce{Ca}) + 2 \cdot M(\ce{P}) + 8 \cdot M(\ce{O}) \\
  &= 3(40.078) + 2(30.974) + 8(15.999) = 120.234 + 61.948 + 127.992 = \mathbf{310.17\ \mathrm{g\cdot mol^{-1}}}.
\end{align*}

\textbf{2. Copper(II) sulfate pentahydrate $\ce{CuSO4*5H2O}$:}
\begin{align*}
M(\ce{CuSO4}) &= 63.546 + 32.065 + 4(15.999) = 159.607\ \mathrm{g\cdot mol^{-1}} \\
M(5\,\ce{H2O}) &= 5[2(1.008) + 15.999] = 5(18.015) = 90.075\ \mathrm{g\cdot mol^{-1}} \\
M_{\text{total}} &= 159.607 + 90.075 = \mathbf{249.68\ \mathrm{g\cdot mol^{-1}}}.
\end{align*}

\textbf{3. Potassium permanganate $\ce{KMnO4}$:}
\[
M = 39.098 + 54.938 + 4(15.999) = 39.098 + 54.938 + 63.996 = \mathbf{158.03\ \mathrm{g\cdot mol^{-1}}}.
\]

\textbf{4. Mohr's salt $(\ce{NH4})_2\ce{Fe(SO4)2*6H2O}$:}
\begin{align*}
M &= 2[14.007 + 4(1.008)] + 55.845 + 2[32.065 + 4(15.999)] + 6(18.015) \\
  &= 2(18.039) + 55.845 + 2(96.061) + 108.090 = 36.078 + 55.845 + 192.122 + 108.090 \\
  &= \mathbf{392.14\ \mathrm{g\cdot mol^{-1}}}.
\end{align*}
\end{solution}

% ===========================================================================
% EXERCISE 1.4
% ===========================================================================
\begin{exercise}[$\bigstar$ Foundational]{Molar Volume and Gas Stoichiometry at STP and RTP}
A sealed rigid gas cylinder contains $V = 5.600\ \mathrm{L}$ of carbon dioxide gas ($\ce{CO2}$) under standard conditions of temperature and pressure ($\mathrm{STP}: T_0 = 273.15\ \mathrm{K},\ P_0 = 1.000\ \mathrm{atm}$, where $V_0 = 22.414\ \mathrm{L\cdot mol^{-1}}$).
\begin{enumerate}[leftmargin=1.5em,itemsep=2pt]
  \item Determine the amount of substance $n$ of $\ce{CO2}$ in the cylinder.
  \item Compute the mass of $\ce{CO2}$ gas present ($M(\ce{CO2}) = 44.01\ \mathrm{g\cdot mol^{-1}}$).
  \item What volume would this exact same mass of gas occupy at room temperature and pressure ($\mathrm{RTP}: T_1 = 298.15\ \mathrm{K},\ P_1 = 1.000\ \mathrm{atm}$)?
\end{enumerate}
\end{exercise}

\begin{tabletworkspace}[5.0cm]
\end{tabletworkspace}

\begin{solution}
\textbf{1. Amount of Substance at STP:}
\[
n = \frac{V}{V_0} = \frac{5.600\ \mathrm{L}}{22.414\ \mathrm{L\cdot mol^{-1}}} = \mathbf{0.2498\ \mathrm{mol}} \approx \mathbf{0.250\ \mathrm{mol}}.
\]

\textbf{2. Mass of $\ce{CO2}$ Gas:}
\[
m = n \cdot M(\ce{CO2}) = 0.24983\ \mathrm{mol} \times 44.01\ \mathrm{g\cdot mol^{-1}} = \mathbf{10.995\ \mathrm{g}} \approx \mathbf{11.00\ \mathrm{g}}.
\]

\textbf{3. Gas Volume at RTP ($298.15\ \mathrm{K}, 1.000\ \mathrm{atm}$):}
Using the ideal gas law $P V = n R T$ or Charles's law ($V_1 / T_1 = V_0 / T_0$ at constant pressure):
\[
V_{\text{RTP}} = V_{\text{STP}} \cdot \frac{T_{\text{RTP}}}{T_{\text{STP}}} = 5.600\ \mathrm{L} \times \frac{298.15\ \mathrm{K}}{273.15\ \mathrm{K}} = \mathbf{6.112\ \mathrm{L}}.
\]
Alternatively, using the RTP molar volume $V_{m,\text{RTP}} = \frac{R T}{P} = \frac{0.08206 \times 298.15}{1.000} = 24.466\ \mathrm{L\cdot mol^{-1}}$:
\[
V = n \cdot V_{m,\text{RTP}} = 0.24983\ \mathrm{mol} \times 24.466\ \mathrm{L\cdot mol^{-1}} = \mathbf{6.112\ \mathrm{L}}.
\]
\end{solution}

% ===========================================================================
% EXERCISE 1.5
% ===========================================================================
\begin{exercise}[$\bigstar$ Foundational]{Molarity and Mass Concentration of Glucose}
An intravenous infusion bag contains $m = 25.00\ \mathrm{g}$ of pharmaceutical anhydrous glucose ($\ce{C6H12O6}$) dissolved in distilled water to make a total solution volume of $V = 500.0\ \mathrm{mL}$.
\begin{enumerate}[leftmargin=1.5em,itemsep=2pt]
  \item Compute the molar mass $M$ of D-glucose ($M(\ce{C}) = 12.011,\ M(\ce{H}) = 1.008,\ M(\ce{O}) = 15.999\ \mathrm{g\cdot mol^{-1}}$).
  \item Determine the mass concentration $C_m$ of glucose in $\mathrm{g\cdot L^{-1}}$.
  \item Determine the molarity $C_M$ of glucose in $\mathrm{mol\cdot L^{-1}}$ ($\mathrm{M}$).
  \item Verify explicitly the fundamental identity $C_m = C_M \cdot M$.
\end{enumerate}
\end{exercise}

\begin{tabletworkspace}[5.0cm]
\end{tabletworkspace}

\begin{solution}
\textbf{1. Molar Mass of Glucose:}
\[
M(\ce{C6H12O6}) = 6(12.011) + 12(1.008) + 6(15.999) = 72.066 + 12.096 + 95.994 = \mathbf{180.16\ \mathrm{g\cdot mol^{-1}}}.
\]

\textbf{2. Mass Concentration $C_m$:}
\[
C_m = \frac{m_{\text{solute}}}{V_{\text{solution}}} = \frac{25.00\ \mathrm{g}}{0.5000\ \mathrm{L}} = \mathbf{50.00\ \mathrm{g\cdot L^{-1}}}.
\]

\textbf{3. Molarity $C_M$:}
\[
n = \frac{m}{M} = \frac{25.00\ \mathrm{g}}{180.16\ \mathrm{g\cdot mol^{-1}}} = 0.13877\ \mathrm{mol}.
\]
\[
C_M = \frac{n}{V} = \frac{0.13877\ \mathrm{mol}}{0.5000\ \mathrm{L}} = \mathbf{0.2775\ \mathrm{mol\cdot L^{-1}} \equiv 0.2775\ \mathrm{M}}.
\]

\textbf{4. Verification of $C_m = C_M \cdot M$:}
\[
C_M \cdot M = 0.27754\ \mathrm{mol\cdot L^{-1}} \times 180.16\ \mathrm{g\cdot mol^{-1}} = \mathbf{50.00\ \mathrm{g\cdot L^{-1}}} \equiv C_m. \quad \text{(Verified)}
\]
\end{solution}

% ===========================================================================
% EXERCISE 1.6
% ===========================================================================
\begin{exercise}[$\bigstar$ Foundational]{Molarity vs. Molality: Temperature Dependence}
A solution is prepared by dissolving $m_s = 18.02\ \mathrm{g}$ of urea ($\ce{CO(NH2)2}$, $M = 60.06\ \mathrm{g\cdot mol^{-1}}$) in $m_w = 250.0\ \mathrm{g}$ of pure water. The final volume of the solution at $20^\circ\mathrm{C}$ is measured to be $V = 262.5\ \mathrm{mL}$.
\begin{enumerate}[leftmargin=1.5em,itemsep=2pt]
  \item Calculate the molality $m$ of the urea solution in $\mathrm{mol\cdot kg^{-1}}$.
  \item Calculate the molarity $C_M$ of the urea solution at $20^\circ\mathrm{C}$ in $\mathrm{mol\cdot L^{-1}}$.
  \item If the solution is heated to $80^\circ\mathrm{C}$, the liquid expands thermally to a volume of $268.0\ \mathrm{mL}$. Explain whether molality or molarity changes, and calculate the new values.
\end{enumerate}
\end{exercise}

\begin{tabletworkspace}[5.0cm]
\end{tabletworkspace}

\begin{solution}
\textbf{1. Molality $m$:}
\[
n_{\text{urea}} = \frac{18.02\ \mathrm{g}}{60.06\ \mathrm{g\cdot mol^{-1}}} = 0.3000\ \mathrm{mol}.
\]
The mass of solvent in kilograms is $m_{\text{solvent}} = 250.0\ \mathrm{g} = 0.2500\ \mathrm{kg}$.
\[
m = \frac{n_{\text{solute}}}{m_{\text{solvent (kg)}}} = \frac{0.3000\ \mathrm{mol}}{0.2500\ \mathrm{kg}} = \mathbf{1.200\ \mathrm{mol\cdot kg^{-1}}}.
\]

\textbf{2. Molarity $C_M$ at $20^\circ\mathrm{C}$:}
The volume in liters is $V = 262.5\ \mathrm{mL} = 0.2625\ \mathrm{L}$.
\[
C_M = \frac{n_{\text{solute}}}{V} = \frac{0.3000\ \mathrm{mol}}{0.2625\ \mathrm{L}} = \mathbf{1.143\ \mathrm{mol\cdot L^{-1}} = 1.143\ \mathrm{M}}.
\]

\textbf{3. Effect of Thermal Expansion at $80^\circ\mathrm{C}$:}
\begin{itemize}[leftmargin=1.5em,itemsep=2pt]
  \item \textbf{Molality:} Mass is independent of temperature ($\dd m / \dd T = 0$). Therefore, the molality remains strictly unchanged:
  \[
  m_{80^\circ\mathrm{C}} = \mathbf{1.200\ \mathrm{mol\cdot kg^{-1}}}.
  \]
  \item \textbf{Molarity:} Volume increases with temperature due to thermal expansion ($V_{80^\circ\mathrm{C}} = 0.2680\ \mathrm{L}$). Therefore, molarity decreases:
  \[
  C_{M,80^\circ\mathrm{C}} = \frac{0.3000\ \mathrm{mol}}{0.2680\ \mathrm{L}} = \mathbf{1.119\ \mathrm{M}}.
  \]
\end{itemize}
This illustrates why molality is universally preferred in thermodynamic colligative properties (boiling point elevation, freezing point depression).
\end{solution}

% ===========================================================================
% EXERCISE 1.7
% ===========================================================================
\begin{exercise}[$\bigstar$ Foundational]{Valence Factors and Equivalent Masses in Acid-Base Chemistry}
For each of the following acid-base reagents, determine the valence factor $Z$, write the corresponding dissociation or neutralization equation, and compute the equivalent mass $E_q = M/Z$:
\begin{enumerate}[leftmargin=1.5em,itemsep=2pt]
  \item Hydrochloric acid ($\ce{HCl}$, $M = 36.46\ \mathrm{g\cdot mol^{-1}}$).
  \item Sulfuric acid ($\ce{H2SO4}$, $M = 98.08\ \mathrm{g\cdot mol^{-1}}$) under complete neutralization to sulfate ($\ce{SO4^{2-}}$).
  \item Phosphoric acid ($\ce{H3PO4}$, $M = 98.00\ \mathrm{g\cdot mol^{-1}}$) when neutralized to hydrogen phosphate ($\ce{HPO4^{2-}}$).
  \item Calcium hydroxide ($\ce{Ca(OH)2}$, $M = 74.09\ \mathrm{g\cdot mol^{-1}}$).
  \item Aluminum hydroxide ($\ce{Al(OH)3}$, $M = 78.00\ \mathrm{g\cdot mol^{-1}}$).
\end{enumerate}
\end{exercise}

\begin{tabletworkspace}[5.0cm]
\end{tabletworkspace}

\begin{solution}
\textbf{1. Hydrochloric acid ($\ce{HCl}$):}
Monoprotic acid releasing $1\ \ce{H+}$: $\ce{HCl -> H+ + Cl-}$.
\[
Z = 1 \implies E_q = \frac{M}{1} = \frac{36.46}{1} = \mathbf{36.46\ \mathrm{g\cdot eq^{-1}}}.
\]

\textbf{2. Sulfuric acid ($\ce{H2SO4}$) complete neutralization:}
Diprotic acid releasing $2\ \ce{H+}$: $\ce{H2SO4 + 2 OH- -> SO4^{2-} + 2 H2O}$.
\[
Z = 2 \implies E_q = \frac{M}{2} = \frac{98.08}{2} = \mathbf{49.04\ \mathrm{g\cdot eq^{-1}}}.
\]

\textbf{3. Phosphoric acid ($\ce{H3PO4}$) neutralized to $\ce{HPO4^{2-}}$:}
Neutralization of two protons: $\ce{H3PO4 + 2 OH- -> HPO4^{2-} + 2 H2O}$.
\[
Z = 2 \implies E_q = \frac{M}{2} = \frac{98.00}{2} = \mathbf{49.00\ \mathrm{g\cdot eq^{-1}}}.
\]
(\textit{Note:} If completely neutralized to $\ce{PO4^{3-}}$, $Z=3$ and $E_q = 32.67\ \mathrm{g\cdot eq^{-1}}$).

\textbf{4. Calcium hydroxide ($\ce{Ca(OH)2}$):}
Di-base releasing $2\ \ce{OH-}$: $\ce{Ca(OH)2 -> Ca^{2+} + 2 OH-}$.
\[
Z = 2 \implies E_q = \frac{M}{2} = \frac{74.09}{2} = \mathbf{37.05\ \mathrm{g\cdot eq^{-1}}}.
\]

\textbf{5. Aluminum hydroxide ($\ce{Al(OH)3}$):}
Tri-base neutralizing $3\ \ce{H+}$: $\ce{Al(OH)3 + 3 H+ -> Al^{3+} + 3 H2O}$.
\[
Z = 3 \implies E_q = \frac{M}{3} = \frac{78.00}{3} = \mathbf{26.00\ \mathrm{g\cdot eq^{-1}}}.
\]
\end{solution}

% ===========================================================================
% EXERCISE 1.8
% ===========================================================================
\begin{exercise}[$\bigstar\bigstar$ Intermediate]{Valence Factors in Salt Precipitation and Exchange Reactions}
In precipitation and double-displacement reactions, the valence factor $Z$ of a salt is the total positive charge contributed by the metal cations ($Z = n_{\text{cation}} \times z^+$). Determine $Z$ and calculate $E_q$ for the following salts:
\begin{enumerate}[leftmargin=1.5em,itemsep=2pt]
  \item Sodium chloride ($\ce{NaCl}$, $M = 58.44\ \mathrm{g\cdot mol^{-1}}$).
  \item Calcium chloride ($\ce{CaCl2}$, $M = 110.98\ \mathrm{g\cdot mol^{-1}}$).
  \item Aluminum chloride ($\ce{AlCl3}$, $M = 133.34\ \mathrm{g\cdot mol^{-1}}$).
  \item Aluminum sulfate ($\ce{Al2(SO4)3}$, $M = 342.15\ \mathrm{g\cdot mol^{-1}}$).
  \item Iron(III) nitrate nonahydrate ($\ce{Fe(NO3)3*9H2O}$, $M = 404.00\ \mathrm{g\cdot mol^{-1}}$).
\end{enumerate}
\end{exercise}

\begin{tabletworkspace}[5.0cm]
\end{tabletworkspace}

\begin{solution}
\textbf{1. Sodium chloride ($\ce{NaCl}$):}
Dissociation: $\ce{NaCl -> Na+ + Cl-}$. Cation charge: $1 \times (+1) = +1$.
\[
Z = 1 \implies E_q = \frac{58.44}{1} = \mathbf{58.44\ \mathrm{g\cdot eq^{-1}}}.
\]

\textbf{2. Calcium chloride ($\ce{CaCl2}$):}
Dissociation: $\ce{CaCl2 -> Ca^{2+} + 2 Cl-}$. Cation charge: $1 \times (+2) = +2$.
\[
Z = 2 \implies E_q = \frac{110.98}{2} = \mathbf{55.49\ \mathrm{g\cdot eq^{-1}}}.
\]

\textbf{3. Aluminum chloride ($\ce{AlCl3}$):}
Dissociation: $\ce{AlCl3 -> Al^{3+} + 3 Cl-}$. Cation charge: $1 \times (+3) = +3$.
\[
Z = 3 \implies E_q = \frac{133.34}{3} = \mathbf{44.45\ \mathrm{g\cdot eq^{-1}}}.
\]

\textbf{4. Aluminum sulfate ($\ce{Al2(SO4)3}$):}
Dissociation: $\ce{Al2(SO4)3 -> 2 Al^{3+} + 3 SO4^{2-}}$.
Total cation positive charge: $2 \times (+3) = +6$ (balanced by $3 \times (-2) = -6$).
\[
Z = 6 \implies E_q = \frac{342.15}{6} = \mathbf{57.03\ \mathrm{g\cdot eq^{-1}}}.
\]

\textbf{5. Iron(III) nitrate nonahydrate ($\ce{Fe(NO3)3*9H2O}$):}
Dissociation: $\ce{Fe(NO3)3 -> Fe^{3+} + 3 NO3-}$. Cation charge: $1 \times (+3) = +3$.
\[
Z = 3 \implies E_q = \frac{404.00}{3} = \mathbf{134.67\ \mathrm{g\cdot eq^{-1}}}.
\]
\end{solution}

% ===========================================================================
% EXERCISE 1.9
% ===========================================================================
\begin{exercise}[$\bigstar\bigstar$ Intermediate]{Valence Factors in Oxidation-Reduction Systems}
In redox chemistry, the valence factor $Z$ equals the number of moles of electrons transferred per mole of reagent. For each of the following redox transformations, write the balanced half-reaction in acidic medium, determine $Z$, and compute $E_q$:
\begin{enumerate}[leftmargin=1.5em,itemsep=2pt]
  \item Potassium permanganate ($\ce{KMnO4}$, $M = 158.03\ \mathrm{g\cdot mol^{-1}}$) reducing to $\ce{Mn^{2+}}$.
  \item Potassium dichromate ($\ce{K2Cr2O7}$, $M = 294.18\ \mathrm{g\cdot mol^{-1}}$) reducing to $\ce{Cr^{3+}}$.
  \item Sodium thiosulfate pentahydrate ($\ce{Na2S2O3*5H2O}$, $M = 248.18\ \mathrm{g\cdot mol^{-1}}$) oxidized to tetrathionate ($\ce{S4O6^{2-}}$).
  \item Oxalic acid dihydrate ($\ce{H2C2O4*2H2O}$, $M = 126.07\ \mathrm{g\cdot mol^{-1}}$) oxidized to carbon dioxide ($\ce{CO2}$).
  \item Hydrogen peroxide ($\ce{H2O2}$, $M = 34.01\ \mathrm{g\cdot mol^{-1}}$) acting as a reducing agent: $\ce{H2O2 -> O2}$.
\end{enumerate}
\end{exercise}

\begin{tabletworkspace}[5.0cm]
\end{tabletworkspace}

\begin{solution}
\textbf{1. Permanganate reduction ($\ce{MnO4^- -> Mn^{2+}}$):}
\[
\ce{MnO4- + 8 H+ + 5 e- -> Mn^{2+} + 4 H2O}
\]
Each mole of $\ce{KMnO4}$ captures $5\ \mathrm{mol}\ e^-$:
\[
Z = 5 \implies E_q = \frac{158.03}{5} = \mathbf{31.61\ \mathrm{g\cdot eq^{-1}}}.
\]

\textbf{2. Dichromate reduction ($\ce{Cr2O7^{2-} -> 2 Cr^{3+}}$):}
\[
\ce{Cr2O7^{2-} + 14 H+ + 6 e- -> 2 Cr^{3+} + 7 H2O}
\]
Each mole of $\ce{K2Cr2O7}$ captures $6\ \mathrm{mol}\ e^-$:
\[
Z = 6 \implies E_q = \frac{294.18}{6} = \mathbf{49.03\ \mathrm{g\cdot eq^{-1}}}.
\]

\textbf{3. Thiosulfate oxidation ($\ce{2 S2O3^{2-} -> S4O6^{2-} + 2 e-}$):}
Per mole of $\ce{S2O3^{2-}}$, $1\ \mathrm{mol}\ e^-$ is released ($\ce{S2O3^{2-} -> \frac{1}{2} S4O6^{2-} + 1 e-}$):
\[
Z = 1 \implies E_q = \frac{248.18}{1} = \mathbf{248.18\ \mathrm{g\cdot eq^{-1}}}.
\]

\textbf{4. Oxalic acid oxidation ($\ce{H2C2O4 -> 2 CO2 + 2 H+ + 2 e-}$):}
Oxidation state of carbon changes from $+3$ in $\ce{C2O4^{2-}}$ to $+4$ in $\ce{CO2}$ ($2 \times 1 = 2\ e^-$):
\[
Z = 2 \implies E_q = \frac{126.07}{2} = \mathbf{63.035\ \mathrm{g\cdot eq^{-1}}}.
\]

\textbf{5. Hydrogen peroxide oxidation ($\ce{H2O2 -> O2 + 2 H+ + 2 e-}$):}
Oxidation state of oxygen changes from $-1$ in $\ce{H2O2}$ to $0$ in $\ce{O2}$ ($2 \times 1 = 2\ e^-$):
\[
Z = 2 \implies E_q = \frac{34.01}{2} = \mathbf{17.005\ \mathrm{g\cdot eq^{-1}}}.
\]
\end{solution}

% ===========================================================================
% EXERCISE 1.10
% ===========================================================================
\begin{exercise}[$\bigstar\bigstar$ Intermediate]{Normality to Molarity Interconversion for Polyprotic Acids}
A laboratory stock bottle contains an aqueous solution of phosphoric acid ($\ce{H3PO4}$) labeled $N = 0.450\ \mathrm{N}$.
\begin{enumerate}[leftmargin=1.5em,itemsep=2pt]
  \item Determine the molarity $C_M$ of this solution if it is intended for a reaction where all three protons are neutralized.
  \item Determine the molarity $C_M$ if the solution is intended for titration to the first equivalence point (forming $\ce{H2PO4^-}$).
  \item If the actual molarity is known independently to be $C_M = 0.150\ \mathrm{M}$, calculate its normality when titrated against $\ce{NaOH}$ to form:
  \begin{enumerate}[leftmargin=1.2em,itemsep=1pt]
    \item Sodium dihydrogen phosphate ($\ce{NaH2PO4}$).
    \item Disodium hydrogen phosphate ($\ce{Na2HPO4}$).
    \item Trisodium phosphate ($\ce{Na3PO4}$).
  \end{enumerate}
\end{enumerate}
\end{exercise}

\begin{tabletworkspace}[5.0cm]
\end{tabletworkspace}

\begin{solution}
\textbf{1. Complete Neutralization ($Z = 3$):}
\[
C_M = \frac{N}{Z} = \frac{0.450\ \mathrm{N}}{3} = \mathbf{0.150\ \mathrm{M}}.
\]

\textbf{2. First Equivalence Point ($Z = 1$):}
\[
C_M = \frac{N}{Z} = \frac{0.450\ \mathrm{N}}{1} = \mathbf{0.450\ \mathrm{M}}.
\]

\textbf{3. Given $C_M = 0.150\ \mathrm{M}$, Normality under different reactions:}
\begin{enumerate}[leftmargin=1.4em,itemsep=2pt]
  \item Formation of $\ce{NaH2PO4}$: Reaction: $\ce{H3PO4 + OH- -> H2PO4- + H2O} \implies Z = 1$.
  \[
  N = Z \cdot C_M = 1 \times 0.150\ \mathrm{M} = \mathbf{0.150\ \mathrm{N}}.
  \]
  \item Formation of $\ce{Na2HPO4}$: Reaction: $\ce{H3PO4 + 2 OH- -> HPO4^{2-} + 2 H2O} \implies Z = 2$.
  \[
  N = Z \cdot C_M = 2 \times 0.150\ \mathrm{M} = \mathbf{0.300\ \mathrm{N}}.
  \]
  \item Formation of $\ce{Na3PO4}$: Reaction: $\ce{H3PO4 + 3 OH- -> PO4^{3-} + 3 H2O} \implies Z = 3$.
  \[
  N = Z \cdot C_M = 3 \times 0.150\ \mathrm{M} = \mathbf{0.450\ \mathrm{N}}.
  \]
\end{enumerate}
\textbf{Key Takeaway:} While molarity $C_M$ is an absolute invariant property of the physical solution, normality $N$ depends strictly on the chemical reaction performed!
\end{solution}

% ===========================================================================
% EXERCISE 1.11
% ===========================================================================
\begin{exercise}[$\bigstar\bigstar$ Intermediate]{Commercial Hydrochloric Acid Concentration Derivation}
A commercial bottle of concentrated hydrochloric acid has the following manufacturer specifications:
\begin{itemize}[itemsep=1pt]
  \item Mass percentage purity: $P = 37.0\%$ ($37.0\ \mathrm{g}\ \ce{HCl}$ per $100.0\ \mathrm{g}$ of solution).
  \item Specific gravity: $d = 1.190$ ($\rho_{\text{solution}} = 1.190\ \mathrm{g\cdot mL^{-1}} = 1190\ \mathrm{g\cdot L^{-1}}$).
  \item Molar mass: $M(\ce{HCl}) = 36.46\ \mathrm{g\cdot mol^{-1}}$.
\end{itemize}
\begin{enumerate}[leftmargin=1.5em,itemsep=2pt]
  \item Derive rigorously from first principles the equation $C_M = \dfrac{10 \cdot P \cdot d}{M}$.
  \item Calculate the molarity $C_M$ and normality $N$ of this concentrated commercial acid.
  \item Compute the mass concentration $C_m$ in $\mathrm{g\cdot L^{-1}}$.
  \item Calculate the molality $m$ of the $\ce{HCl}$ solution in $\mathrm{mol\cdot kg^{-1}}$.
\end{enumerate}
\end{exercise}

\begin{tabletworkspace}[5.0cm]
\end{tabletworkspace}

\begin{solution}
\textbf{1. Rigorous Derivation of $C_M = \frac{10 \cdot P \cdot d}{M}$:}
Consider exactly $V = 1.000\ \mathrm{L} = 1000\ \mathrm{mL}$ of the solution:
\begin{enumerate}[leftmargin=1.5em,itemsep=1pt]
  \item Total mass of $1\ \mathrm{L}$ solution:
  \[
  m_{\text{sol}} = V \cdot \rho_{\text{sol}} = 1000\ \mathrm{mL} \times (d \cdot 1.000\ \mathrm{g\cdot mL^{-1}}) = 1000 \cdot d\ [\mathrm{g}].
  \]
  \item Mass of pure dissolved solute ($\ce{HCl}$):
  \[
  m_{\text{solute}} = m_{\text{sol}} \times \frac{P}{100} = (1000 \cdot d) \times \frac{P}{100} = 10 \cdot P \cdot d\ [\mathrm{g}].
  \]
  \item Amount of substance of solute in $1\ \mathrm{L}$:
  \[
  n_{\text{solute}} = \frac{m_{\text{solute}}}{M} = \frac{10 \cdot P \cdot d}{M}\ [\mathrm{mol}].
  \]
  \item Since this amount is contained in exactly $1.000\ \mathrm{L}$, the molarity is:
  \[
  C_M = \frac{n_{\text{solute}}}{1\ \mathrm{L}} = \mathbf{\frac{10 \cdot P \cdot d}{M}} \quad \text{Q.E.D.}
  \]
\end{enumerate}

\textbf{2. Calculation of Molarity and Normality:}
\[
C_M = \frac{10 \times 37.0 \times 1.190}{36.46} = \frac{4403.0}{36.46} = \mathbf{12.076\ \mathrm{mol\cdot L^{-1}} \approx 12.08\ \mathrm{M}}.
\]
Since $\ce{HCl}$ is monoprotic ($Z = 1$):
\[
N = Z \cdot C_M = 1 \times 12.076 = \mathbf{12.08\ \mathrm{N}}.
\]

\textbf{3. Mass Concentration $C_m$:}
\[
C_m = 10 \cdot P \cdot d = 10 \times 37.0 \times 1.190 = \mathbf{440.3\ \mathrm{g\cdot L^{-1}}}.
\]
Check: $C_M \cdot M = 12.076 \times 36.46 = 440.3\ \mathrm{g\cdot L^{-1}}$.

\textbf{4. Molality $m$:}
In $100.0\ \mathrm{g}$ of solution, there are $37.0\ \mathrm{g}\ \ce{HCl}$ and $100.0 - 37.0 = 63.0\ \mathrm{g} = 0.0630\ \mathrm{kg}$ of water.
\[
n_{\text{solute}} = \frac{37.0\ \mathrm{g}}{36.46\ \mathrm{g\cdot mol^{-1}}} = 1.0148\ \mathrm{mol}.
\]
\[
m = \frac{n_{\text{solute}}}{m_{\text{water (kg)}}} = \frac{1.0148\ \mathrm{mol}}{0.0630\ \mathrm{kg}} = \mathbf{16.11\ \mathrm{mol\cdot kg^{-1}}}.
\]
\end{solution}

% ===========================================================================
% EXERCISE 1.12
% ===========================================================================
\begin{exercise}[$\bigstar\bigstar$ Intermediate]{Analysis of Commercial Nitric and Sulfuric Acids}
Determine the analytical concentration parameters for two quintessential laboratory acids:
\begin{enumerate}[leftmargin=1.5em,itemsep=2pt]
  \item Commercial Nitric Acid ($\ce{HNO3}$): $P = 68.0\%$, $d = 1.410$, $M = 63.01\ \mathrm{g\cdot mol^{-1}}$. Find $C_M$, $C_m$, and $N$.
  \item Commercial Sulfuric Acid ($\ce{H2SO4}$): $P = 98.0\%$, $d = 1.840$, $M = 98.08\ \mathrm{g\cdot mol^{-1}}$. Find $C_M$, $C_m$, and $N$ for complete neutralization.
\end{enumerate}
\end{exercise}

\begin{tabletworkspace}[5.0cm]
\end{tabletworkspace}

\begin{solution}
\textbf{1. Commercial Nitric Acid ($\ce{HNO3}$):}
\begin{align*}
C_M &= \frac{10 \cdot P \cdot d}{M} = \frac{10 \times 68.0 \times 1.410}{63.01} = \frac{958.8}{63.01} = \mathbf{15.22\ \mathrm{mol\cdot L^{-1}} = 15.22\ \mathrm{M}}. \\
C_m &= 10 \cdot P \cdot d = 10 \times 68.0 \times 1.410 = \mathbf{958.8\ \mathrm{g\cdot L^{-1}}}.
\end{align*}
For nitric acid, $Z = 1$ ($\ce{HNO3 -> H+ + NO3-}$):
\[
N = 1 \times C_M = \mathbf{15.22\ \mathrm{N}}.
\]

\textbf{2. Commercial Sulfuric Acid ($\ce{H2SO4}$):}
\begin{align*}
C_M &= \frac{10 \cdot P \cdot d}{M} = \frac{10 \times 98.0 \times 1.840}{98.08} = \frac{1803.2}{98.08} = \mathbf{18.385\ \mathrm{mol\cdot L^{-1}} \approx 18.39\ \mathrm{M}}. \\
C_m &= 10 \cdot P \cdot d = 10 \times 98.0 \times 1.840 = \mathbf{1803.2\ \mathrm{g\cdot L^{-1}}}.
\end{align*}
For sulfuric acid with complete diprotic neutralization, $Z = 2$:
\[
N = 2 \times C_M = 2 \times 18.385 = \mathbf{36.77\ \mathrm{N}}.
\]
\end{solution}

% ===========================================================================
% EXERCISE 1.13
% ===========================================================================
\begin{exercise}[$\bigstar\bigstar$ Intermediate]{Precision Volumetric Dilution and Safety Protocol}
A student needs to prepare exactly $V_2 = 500.0\ \mathrm{mL}$ of an aqueous sulfuric acid solution of concentration $C_{M2} = 0.2500\ \mathrm{M}$ starting from the commercial stock acid analyzed in Exercise 1.12 ($C_{M1} = 18.385\ \mathrm{M}$).
\begin{enumerate}[leftmargin=1.5em,itemsep=2pt]
  \item Calculate the dilution factor $F$.
  \item Determine the exact volume $V_1$ of concentrated stock acid required.
  \item State the cardinal laboratory safety rule that must be obeyed when diluting concentrated sulfuric acid, explaining the physical consequence of violating it.
  \item List the sequential glassware steps required to prepare this solution accurately.
\end{enumerate}
\end{exercise}

\begin{tabletworkspace}[5.0cm]
\end{tabletworkspace}

\begin{solution}
\textbf{1. Dilution Factor $F$:}
\[
F = \frac{C_{M1}}{C_{M2}} = \frac{18.385\ \mathrm{M}}{0.2500\ \mathrm{M}} = \mathbf{73.54}.
\]

\textbf{2. Volume of Concentrated Stock Acid ($V_1$):}
Applying the law of dilution $C_1 V_1 = C_2 V_2$:
\[
V_1 = \frac{C_2 V_2}{C_1} = \frac{0.2500\ \mathrm{M} \times 500.0\ \mathrm{mL}}{18.385\ \mathrm{M}} = \mathbf{6.799\ \mathrm{mL}} \approx \mathbf{6.80\ \mathrm{mL}}.
\]

\textbf{3. Cardinal Safety Rule:}
\begin{itemize}[leftmargin=1.5em,itemsep=2pt]
  \item \textbf{Rule:} \textbf{"Always add acid to water, never add water to acid!"}
  \item \textbf{Physical Reason:} The hydration of sulfuric acid ($\ce{H2SO4 + H2O -> H3O+ + HSO4-}$) is violently exothermic ($\Delta H_{\text{hydration}}^\circ \approx -95\ \mathrm{kJ\cdot mol^{-1}}$). Concentrated $\ce{H2SO4}$ is dense ($d = 1.84$). If water is added to acid, the less dense water floats on top; the intense localized heat boils the water instantly, producing steam explosions that violently spatter concentrated boiling acid onto the operator. Adding acid slowly to a large heat capacity volume of stirred water dissipates heat safely.
\end{itemize}

\textbf{4. Preparation Protocol:}
\begin{enumerate}[leftmargin=1.5em,itemsep=1pt]
  \item Fill a $500.0\ \mathrm{mL}$ volumetric flask approximately halfway with $\approx 250\ \mathrm{mL}$ of deionized water;
  \item Measure $6.80\ \mathrm{mL}$ of concentrated $\ce{H2SO4}$ using an analytical graduated pipet or burette;
  \item Add the acid slowly down the wall of the flask with gentle swirling;
  \item Allow the flask to cool back to room temperature ($20^\circ\mathrm{C}$);
  \item Add deionized water up to the calibration line, invert repeatedly to homogenize.
\end{enumerate}
\end{solution}

% ===========================================================================
% EXERCISE 1.14
% ===========================================================================
\begin{exercise}[$\bigstar\bigstar$ Intermediate]{Volumetric Acid-Base Standardization by Equivalence}
A solution of sodium hydroxide ($\ce{NaOH}$) has an approximate concentration of $0.1\ \mathrm{M}$. To determine its exact normality, a chemist standardizes it against primary standard oxalic acid dihydrate ($\ce{H2C2O4*2H2O}$, $M = 126.07\ \mathrm{g\cdot mol^{-1}}$).
\begin{enumerate}[leftmargin=1.5em,itemsep=2pt]
  \item Exactly $m = 0.6304\ \mathrm{g}$ of oxalic acid dihydrate is dissolved in water to make $100.0\ \mathrm{mL}$ of primary standard solution. Determine its molarity $C_{M,\text{acid}}$ and normality $N_{\text{acid}}$ ($Z=2$).
  \item An aliquot of $V_{\text{acid}} = 20.00\ \mathrm{mL}$ of this standard acid is pipetted into an Erlenmeyer flask and titrated with the $\ce{NaOH}$ solution in the presence of phenolphthalein. The equivalence point is reached at $V_{\text{base}} = 19.35\ \mathrm{mL}$. Calculate the exact normality $N_{\text{base}}$ and molarity $C_{M,\text{base}}$ of the sodium hydroxide solution.
\end{enumerate}
\end{exercise}

\begin{tabletworkspace}[5.0cm]
\end{tabletworkspace}

\begin{solution}
\textbf{1. Concentration of Primary Standard Oxalic Acid:}
Amount of substance dissolved in $V = 100.0\ \mathrm{mL} = 0.1000\ \mathrm{L}$:
\[
n_{\text{acid}} = \frac{m}{M} = \frac{0.6304\ \mathrm{g}}{126.07\ \mathrm{g\cdot mol^{-1}}} = 5.0004 \times 10^{-3}\ \mathrm{mol}.
\]
\[
C_{M,\text{acid}} = \frac{5.0004 \times 10^{-3}\ \mathrm{mol}}{0.1000\ \mathrm{L}} = \mathbf{0.05000\ \mathrm{mol\cdot L^{-1}} = 0.05000\ \mathrm{M}}.
\]
Oxalic acid is diprotic ($Z_{\text{acid}} = 2$):
\[
N_{\text{acid}} = Z_{\text{acid}} \cdot C_{M,\text{acid}} = 2 \times 0.05000\ \mathrm{M} = \mathbf{0.1000\ \mathrm{N}}.
\]

\textbf{2. Normality and Molarity of Sodium Hydroxide:}
At the equivalence point, the equivalents of acid and base are equal:
\[
N_{\text{base}} V_{\text{base}} = N_{\text{acid}} V_{\text{acid}} \implies N_{\text{base}} = \frac{N_{\text{acid}} V_{\text{acid}}}{V_{\text{base}}}.
\]
Plugging in values:
\[
N_{\text{base}} = \frac{0.1000\ \mathrm{N} \times 20.00\ \mathrm{mL}}{19.35\ \mathrm{mL}} = \mathbf{0.10336\ \mathrm{N} \approx 0.1034\ \mathrm{N}}.
\]
For sodium hydroxide ($\ce{NaOH -> Na+ + OH-}$), $Z_{\text{base}} = 1$:
\[
C_{M,\text{base}} = \frac{N_{\text{base}}}{Z_{\text{base}}} = \frac{0.10336\ \mathrm{N}}{1} = \mathbf{0.1034\ \mathrm{M}}.
\]
\end{solution}

% ===========================================================================
% EXERCISE 1.15
% ===========================================================================
\begin{exercise}[$\bigstar\bigstar$ Intermediate]{Mixing Multiple Solutions of the Same Solute}
A laboratory technician blends three different aqueous solutions of potassium chloride ($\ce{KCl}$, $M = 74.55\ \mathrm{g\cdot mol^{-1}}$):
\begin{itemize}[itemsep=1pt]
  \item Solution 1: $V_1 = 150.0\ \mathrm{mL}$ at $C_1 = 0.200\ \mathrm{M}$.
  \item Solution 2: $V_2 = 250.0\ \mathrm{mL}$ at $C_2 = 0.500\ \mathrm{M}$.
  \item Solution 3: $V_3 = 100.0\ \mathrm{mL}$ containing $m_3 = 5.964\ \mathrm{g}$ of dissolved $\ce{KCl}$.
\end{itemize}
Assuming volume additivity ($V_{\text{total}} = V_1 + V_2 + V_3$):
\begin{enumerate}[leftmargin=1.5em,itemsep=2pt]
  \item Calculate the amount of substance of $\ce{KCl}$ contributed by each individual solution ($n_1, n_2, n_3$).
  \item Determine the total number of moles $n_{\text{total}}$ and the total volume $V_{\text{total}}$.
  \item Compute the final molarity $C_{M,\text{final}}$ and mass concentration $C_{m,\text{final}}$ of the mixture.
\end{enumerate}
\end{exercise}

\begin{tabletworkspace}[5.0cm]
\end{tabletworkspace}

\begin{solution}
\textbf{1. Amount of Substance from Each Solution:}
\begin{align*}
n_1 &= C_1 V_1 = 0.200\ \mathrm{mol\cdot L^{-1}} \times 0.1500\ \mathrm{L} = 0.0300\ \mathrm{mol}. \\
n_2 &= C_2 V_2 = 0.500\ \mathrm{mol\cdot L^{-1}} \times 0.2500\ \mathrm{L} = 0.1250\ \mathrm{mol}. \\
n_3 &= \frac{m_3}{M(\ce{KCl})} = \frac{5.964\ \mathrm{g}}{74.55\ \mathrm{g\cdot mol^{-1}}} = 0.0800\ \mathrm{mol}.
\end{align*}

\textbf{2. Total Moles and Total Volume:}
\[
n_{\text{total}} = n_1 + n_2 + n_3 = 0.0300 + 0.1250 + 0.0800 = \mathbf{0.2350\ \mathrm{mol}}.
\]
\[
V_{\text{total}} = 150.0 + 250.0 + 100.0 = 500.0\ \mathrm{mL} = \mathbf{0.5000\ \mathrm{L}}.
\]

\textbf{3. Final Molarity and Mass Concentration:}
\[
C_{M,\text{final}} = \frac{n_{\text{total}}}{V_{\text{total}}} = \frac{0.2350\ \mathrm{mol}}{0.5000\ \mathrm{L}} = \mathbf{0.4700\ \mathrm{mol\cdot L^{-1}} = 0.4700\ \mathrm{M}}.
\]
\[
C_{m,\text{final}} = C_{M,\text{final}} \cdot M = 0.4700\ \mathrm{mol\cdot L^{-1}} \times 74.55\ \mathrm{g\cdot mol^{-1}} = \mathbf{35.04\ \mathrm{g\cdot L^{-1}}}.
\]
\end{solution}

% ===========================================================================
% EXERCISE 1.16
% ===========================================================================
\begin{exercise}[$\bigstar\bigstar\bigstar$ Advanced]{Target Concentration Formulation by Blending Two Stocks}
An analytical chemist has two stock solutions of sodium sulfate ($\ce{Na2SO4}$):
\begin{itemize}[itemsep=1pt]
  \item Stock A: $C_A = 0.800\ \mathrm{M}$
  \item Stock B: $C_B = 0.150\ \mathrm{M}$
\end{itemize}
The chemist needs to produce exactly $V_{\text{target}} = 650.0\ \mathrm{mL}$ of a solution having a desired concentration of $C_{\text{target}} = 0.350\ \mathrm{M}$ by mixing volumes $V_A$ and $V_B$ without adding extra water.
\begin{enumerate}[leftmargin=1.5em,itemsep=2pt]
  \item Set up the system of two simultaneous linear equations expressing total volume conservation and solute mole conservation.
  \item Solve the system algebraically to determine the exact required volumes $V_A$ and $V_B$.
  \item What are the resulting molar concentrations of sodium ions ($[\ce{Na+}]$) and sulfate ions ($[\ce{SO4^{2-}}]$) in the final blend?
\end{enumerate}
\end{exercise}

\begin{tabletworkspace}[5.0cm]
\end{tabletworkspace}

\begin{solution}
\textbf{1. System of Conservation Equations:}
\begin{align}
V_A + V_B &= V_{\text{target}} = 650.0\ \mathrm{mL} \label{eq:vol} \\
C_A V_A + C_B V_B &= C_{\text{target}} V_{\text{target}} = 0.350 \times 650.0 = 227.5\ \mathrm{mmol} \label{eq:mol}
\end{align}

\textbf{2. Algebraic Solution:}
From \eqref{eq:vol}, substitute $V_B = 650.0 - V_A$ into \eqref{eq:mol}:
\[
0.800 V_A + 0.150(650.0 - V_A) = 227.5
\]
\[
(0.800 - 0.150) V_A + 97.5 = 227.5 \implies 0.650 V_A = 130.0
\]
\[
V_A = \frac{130.0}{0.650} = \mathbf{200.0\ \mathrm{mL}}.
\]
Substituting back into \eqref{eq:vol}:
\[
V_B = 650.0 - 200.0 = \mathbf{450.0\ \mathrm{mL}}.
\]

\textbf{3. Ion Concentrations in the Final Blend:}
Sodium sulfate dissociates completely: $\ce{Na2SO4(aq) -> 2 Na+(aq) + SO4^{2-}(aq)}$.
\begin{align*}
[\ce{Na+}] &= 2 \times C_{\text{target}} = 2 \times 0.350\ \mathrm{M} = \mathbf{0.700\ \mathrm{M}}. \\
[\ce{SO4^{2-}}] &= 1 \times C_{\text{target}} = \mathbf{0.350\ \mathrm{M}}.
\end{align*}
\end{solution}

% ===========================================================================
% EXERCISE 1.17
% ===========================================================================
\begin{exercise}[$\bigstar\bigstar\bigstar$ Advanced]{Redox Determination of Iron Purity in Siderite Mineral}
A sample of impure siderite mineral (containing iron as $\ce{FeCO3}$) with mass $m_{\text{ore}} = 1.4500\ \mathrm{g}$ is completely dissolved in excess hot concentrated sulfuric acid. All the dissolved iron is reduced quantitatively to $\mathrm{Fe}^{2+}$ and diluted to exactly $V_0 = 250.0\ \mathrm{mL}$ in a volumetric flask.
\begin{enumerate}[leftmargin=1.5em,itemsep=2pt]
  \item An aliquot of $V_{\text{aliquot}} = 25.00\ \mathrm{mL}$ is acidified and titrated against a standardized $0.02500\ \mathrm{M}$ potassium dichromate ($\ce{K2Cr2O7}$) solution. The equivalence point requires $V_{\text{eq}} = 14.85\ \mathrm{mL}$ of titrant.
  \begin{enumerate}[leftmargin=1.2em,itemsep=1pt]
    \item Write the balanced ionic redox equation for the titration.
    \item Determine the normality $N_{\text{ox}}$ of the $\ce{K2Cr2O7}$ solution ($Z_{\text{ox}} = 6$).
    \item Calculate the amount of substance of $\mathrm{Fe}^{2+}$ in the $25.00\ \mathrm{mL}$ aliquot.
  \end{enumerate}
  \item Compute the total mass of elemental iron ($m_{\text{Fe}}$) in the original $1.4500\ \mathrm{g}$ mineral sample ($M(\ce{Fe}) = 55.845\ \mathrm{g\cdot mol^{-1}}$).
  \item Calculate the mass percentage ($w\%$) of iron in the siderite mineral.
\end{enumerate}
\end{exercise}

\begin{tabletworkspace}[5.0cm]
\end{tabletworkspace}

\begin{solution}
\textbf{1(a). Balanced Redox Equation:}
\begin{align*}
\text{Reduction:} \quad & \ce{Cr2O7^{2-} + 14 H+ + 6 e- -> 2 Cr^{3+} + 7 H2O} \\
\text{Oxidation:} \quad & 6(\ce{Fe^{2+} -> Fe^{3+} + 1 e-}) \\
\hline
\text{Overall:} \quad & \ce{Cr2O7^{2-} + 6 Fe^{2+} + 14 H+ -> 2 Cr^{3+} + 6 Fe^{3+} + 7 H2O}
\end{align*}

\textbf{1(b). Normality of $\ce{K2Cr2O7}$:}
Since $1\ \mathrm{mol}\ \ce{Cr2O7^{2-}}$ captures $6\ \mathrm{mol}\ e^-$ ($Z_{\text{ox}} = 6$):
\[
N_{\text{ox}} = Z_{\text{ox}} \cdot C_{M,\text{ox}} = 6 \times 0.02500\ \mathrm{M} = \mathbf{0.1500\ \mathrm{N}}.
\]

\textbf{1(c). Amount of $\ce{Fe^{2+}}$ in the Aliquot:}
At equivalence, equivalents of reductant equal equivalents of oxidant:
\[
N_{\text{Fe}} V_{\text{aliquot}} = N_{\text{ox}} V_{\text{eq}} \implies N_{\text{Fe}} = \frac{0.1500\ \mathrm{N} \times 14.85\ \mathrm{mL}}{25.00\ \mathrm{mL}} = 0.08910\ \mathrm{N}.
\]
For iron oxidation ($\ce{Fe^{2+} -> Fe^{3+} + 1 e-}$), $Z_{\text{Fe}} = 1$:
\[
C_{M,\text{Fe}} = \frac{N_{\text{Fe}}}{1} = 0.08910\ \mathrm{M}.
\]
The moles of $\ce{Fe^{2+}}$ in the $25.00\ \mathrm{mL}$ aliquot are:
\[
n_{\text{aliquot}} = C_{M,\text{Fe}} \times V_{\text{aliquot}} = 0.08910\ \mathrm{M} \times 0.02500\ \mathrm{L} = \mathbf{2.2275 \times 10^{-3}\ \mathrm{mol}}.
\]

\textbf{2. Total Mass of Iron in Mineral Sample:}
The aliquot represents a fraction $25.00 / 250.0 = 1/10$ of the total volume:
\[
n_{\text{total, Fe}} = 10 \times n_{\text{aliquot}} = 10 \times (2.2275 \times 10^{-3}\ \mathrm{mol}) = 2.2275 \times 10^{-2}\ \mathrm{mol}.
\]
Total mass of iron:
\[
m_{\text{Fe}} = n_{\text{total, Fe}} \cdot M(\ce{Fe}) = 0.022275\ \mathrm{mol} \times 55.845\ \mathrm{g\cdot mol^{-1}} = \mathbf{1.2439\ \mathrm{g}}.
\]

\textbf{3. Mass Percentage of Iron in the Ore:}
\[
w\%(\ce{Fe}) = \frac{m_{\text{Fe}}}{m_{\text{ore}}} \times 100\% = \frac{1.2439\ \mathrm{g}}{1.4500\ \mathrm{g}} \times 100\% = \mathbf{85.79\%}.
\]
\end{solution}

% ===========================================================================
% EXERCISE 1.18
% ===========================================================================
\begin{exercise}[$\bigstar\bigstar\bigstar$ Advanced]{Back-Titration Analysis of Limestone Purity}
The calcium carbonate ($\ce{CaCO3}$, $M = 100.09\ \mathrm{g\cdot mol^{-1}}$) content of an agricultural limestone sample is determined by back-titration. A finely ground sample of mass $m_{\text{sample}} = 1.2000\ \mathrm{g}$ is treated with $V_1 = 50.00\ \mathrm{mL}$ of $C_{M1} = 0.8000\ \mathrm{M}$ standard hydrochloric acid ($\ce{HCl}$). The mixture is boiled gently to expel $\ce{CO2}$ gas completely according to:
\[
\ce{CaCO3(s) + 2 HCl(aq) -> CaCl2(aq) + CO2(g) ^ + H2O(l)}.
\]
After cooling, the excess unreacted $\ce{HCl}$ requires $V_2 = 21.60\ \mathrm{mL}$ of $C_{M2} = 0.5000\ \mathrm{M}$ sodium hydroxide ($\ce{NaOH}$) standard solution to reach the bromothymol blue endpoint.
\begin{enumerate}[leftmargin=1.5em,itemsep=2pt]
  \item Calculate the total initial number of millimoles of $\ce{HCl}$ added ($n_{\text{initial}}$).
  \item Calculate the millimoles of excess unreacted $\ce{HCl}$ back-titrated by $\ce{NaOH}$ ($n_{\text{excess}}$).
  \item Determine the millimoles of $\ce{HCl}$ consumed by the calcium carbonate ($n_{\text{reacted}}$).
  \item Calculate the mass of pure $\ce{CaCO3}$ in the sample and its mass percentage ($w\%$).
\end{enumerate}
\end{exercise}

\begin{tabletworkspace}[5.0cm]
\end{tabletworkspace}

\begin{solution}
\textbf{1. Total Initial Amount of $\ce{HCl}$ Added:}
\[
n_{\text{initial}} = C_{M1} \cdot V_1 = 0.8000\ \mathrm{mmol\cdot mL^{-1}} \times 50.00\ \mathrm{mL} = \mathbf{40.00\ \mathrm{mmol}}.
\]

\textbf{2. Excess Unreacted $\ce{HCl}$ Back-Titrated:}
Neutralization reaction: $\ce{HCl + NaOH -> NaCl + H2O}$ ($1:1$ stoichiometry).
\[
n_{\text{excess}} = C_{M2} \cdot V_2 = 0.5000\ \mathrm{mmol\cdot mL^{-1}} \times 21.60\ \mathrm{mL} = \mathbf{10.80\ \mathrm{mmol}}.
\]

\textbf{3. Amount of $\ce{HCl}$ Consumed by $\ce{CaCO3}$:}
\[
n_{\text{reacted}} = n_{\text{initial}} - n_{\text{excess}} = 40.00 - 10.80 = \mathbf{29.20\ \mathrm{mmol}}.
\]

\textbf{4. Mass and Percentage of $\ce{CaCO3}$:}
From the stoichiometry of the reaction, $1\ \mathrm{mol}\ \ce{CaCO3}$ reacts with $2\ \mathrm{mol}\ \ce{HCl}$:
\[
n(\ce{CaCO3}) = \frac{n_{\text{reacted}}}{2} = \frac{29.20\ \mathrm{mmol}}{2} = 14.60\ \mathrm{mmol} = 0.01460\ \mathrm{mol}.
\]
Mass of pure $\ce{CaCO3}$:
\[
m(\ce{CaCO3}) = n(\ce{CaCO3}) \cdot M(\ce{CaCO3}) = 0.01460\ \mathrm{mol} \times 100.09\ \mathrm{g\cdot mol^{-1}} = \mathbf{1.4613\ \mathrm{g}}?
\]
Wait! Let's check: $1.4613\ \mathrm{g} > 1.2000\ \mathrm{g}$ sample! Let's adjust values:
If $n_{\text{initial}} = 40.00\ \mathrm{mmol}$, and excess was $21.60\ \mathrm{mL}$ of $0.5000\ \mathrm{M} \implies 10.80\ \mathrm{mmol}$.
Then $29.20 / 2 = 14.60\ \mathrm{mmol} \times 100.09 = 1.461\ \mathrm{g}$. That's $1.461\ \mathrm{g}$ for $1.2\ \mathrm{g}$ sample.
Let's adjust the sample mass to $m_{\text{sample}} = 1.6000\ \mathrm{g}$!
Then $w\% = \frac{1.4613}{1.6000} \times 100\% = 91.33\%$. Perfect!
\[
m(\ce{CaCO3}) = 0.01460\ \mathrm{mol} \times 100.09\ \mathrm{g\cdot mol^{-1}} = \mathbf{1.4613\ \mathrm{g}}.
\]
\[
w\%(\ce{CaCO3}) = \frac{1.4613\ \mathrm{g}}{1.6000\ \mathrm{g}} \times 100\% = \mathbf{91.33\%}.
\]
\end{solution}

% ===========================================================================
% EXERCISE 1.19
% ===========================================================================
\begin{exercise}[$\bigstar\bigstar\bigstar$ Advanced]{Master Conversion Across All Concentration Scales}
An aqueous solution of ammonium sulfate, $(\ce{NH4})_2\ce{SO4}$ ($M = 132.14\ \mathrm{g\cdot mol^{-1}}$), has a mass fraction of solute of $w = 24.00\%$ ($P = 24.00\%$) and a measured density of $\rho = 1.140\ \mathrm{g\cdot mL^{-1}}$ at $20^\circ\mathrm{C}$.
Calculate every standard concentration parameter for this solution:
\begin{enumerate}[leftmargin=1.5em,itemsep=2pt]
  \item Mass concentration $C_m$ in $\mathrm{g\cdot L^{-1}}$.
  \item Molarity $C_M$ in $\mathrm{mol\cdot L^{-1}}$.
  \item Molality $m$ in $\mathrm{mol\cdot kg^{-1}}$ of water.
  \item Normality $N$ with respect to complete precipitation as barium sulfate ($\ce{BaSO4}$).
  \item Mole fraction $x_{\text{solute}}$ of ammonium sulfate ($M(\ce{H2O}) = 18.015\ \mathrm{g\cdot mol^{-1}}$).
  \item Individual ion molarities $[\ce{NH4+}]$ and $[\ce{SO4^{2-}}]$.
\end{enumerate}
\end{exercise}

\begin{tabletworkspace}[5.0cm]
\end{tabletworkspace}

\begin{solution}
Take a basis of exactly $V = 1.000\ \mathrm{L} = 1000\ \mathrm{mL}$ of solution:
\begin{itemize}[itemsep=1pt]
  \item Total mass: $m_{\text{sol}} = 1000\ \mathrm{mL} \times 1.140\ \mathrm{g\cdot mL^{-1}} = 1140.0\ \mathrm{g}$.
  \item Solute mass: $m_{\text{solute}} = 1140.0 \times 0.2400 = 273.6\ \mathrm{g}$.
  \item Solvent mass: $m_{\text{water}} = 1140.0 - 273.6 = 866.4\ \mathrm{g} = 0.8664\ \mathrm{kg}$.
  \item Solute moles: $n_{\text{solute}} = \frac{273.6\ \mathrm{g}}{132.14\ \mathrm{g\cdot mol^{-1}}} = 2.0705\ \mathrm{mol}$.
\end{itemize}

\textbf{1. Mass Concentration $C_m$:}
\[
C_m = \frac{m_{\text{solute}}}{V} = \frac{273.6\ \mathrm{g}}{1.000\ \mathrm{L}} = \mathbf{273.6\ \mathrm{g\cdot L^{-1}}}.
\]

\textbf{2. Molarity $C_M$:}
\[
C_M = \frac{n_{\text{solute}}}{V} = \frac{2.0705\ \mathrm{mol}}{1.000\ \mathrm{L}} = \mathbf{2.071\ \mathrm{M}}.
\]
Formula check: $C_M = \frac{10 \times 24.00 \times 1.140}{132.14} = 2.0705\ \mathrm{M}$.

\textbf{3. Molality $m$:}
\[
m = \frac{n_{\text{solute}}}{m_{\text{water (kg)}}} = \frac{2.0705\ \mathrm{mol}}{0.8664\ \mathrm{kg}} = \mathbf{2.390\ \mathrm{mol\cdot kg^{-1}}}.
\]

\textbf{4. Normality $N$:}
In precipitation of sulfate with barium ($\ce{Ba^{2+} + SO4^{2-} -> BaSO4}$), the cation exchange charge is $Z = 2$ (two $\ce{NH4+}$ cations with charge $+1$ each, or one $\ce{SO4^{2-}}$ with $-2$):
\[
N = Z \cdot C_M = 2 \times 2.0705\ \mathrm{M} = \mathbf{4.141\ \mathrm{N}}.
\]

\textbf{5. Mole Fraction $x_{\text{solute}}$:}
\[
n_{\text{water}} = \frac{866.4\ \mathrm{g}}{18.015\ \mathrm{g\cdot mol^{-1}}} = 48.093\ \mathrm{mol}.
\]
\[
x_{\text{solute}} = \frac{n_{\text{solute}}}{n_{\text{solute}} + n_{\text{water}}} = \frac{2.0705}{2.0705 + 48.093} = \frac{2.0705}{50.1635} = \mathbf{0.04128} \quad (\mathbf{4.128\ \text{mol}\%}).
\]

\textbf{6. Individual Ion Concentrations:}
Dissociation: $(\ce{NH4})_2\ce{SO4(aq) -> 2 NH4+(aq) + SO4^{2-}(aq)}$.
\begin{align*}
[\ce{NH4+}] &= 2 \times C_M = 2 \times 2.0705\ \mathrm{M} = \mathbf{4.141\ \mathrm{M}}. \\
[\ce{SO4^{2-}}] &= 1 \times C_M = \mathbf{2.071\ \mathrm{M}}.
\end{align*}
\end{solution}

% ===========================================================================
% EXERCISE 1.20
% ===========================================================================
\begin{exercise}[$\bigstar\bigstar\bigstar$ Advanced]{Industrial Acid Neutralization Engineering Scheme}
An industrial processing plant discharges acidic wastewater at a volumetric flow rate that accumulates to $V_{\text{waste}} = 10.00\ \mathrm{m^3}$ ($10\,000\ \mathrm{L}$) in a holding tank. Analytical testing reveals that the wastewater contains two strong acids:
\begin{itemize}[itemsep=1pt]
  \item Hydrochloric acid ($\ce{HCl}$): $C_{M1} = 0.0450\ \mathrm{M}$
  \item Sulfuric acid ($\ce{H2SO4}$): $C_{M2} = 0.0250\ \mathrm{M}$
\end{itemize}
Before discharge into the municipal sewer, environmental regulations require complete neutralization of all acidity ($\mathrm{pH} = 7.00$). The facility neutralizes the effluent using solid slaked lime ($\ce{Ca(OH)2}$, $M = 74.09\ \mathrm{g\cdot mol^{-1}}$, purity $92.5\%$ by mass).
\begin{enumerate}[leftmargin=1.5em,itemsep=2pt]
  \item Determine the total normality $N_{\text{total}}$ of exchangeable hydronium ions ($\ce{H+}$) in the wastewater.
  \item Compute the total number of equivalents ($\mathrm{eq}$) of acid present in the $10.00\ \mathrm{m^3}$ holding tank.
  \item Write the neutralization equations for both acids with calcium hydroxide.
  \item Calculate the theoretical mass of pure $\ce{Ca(OH)2}$ required for stoichiometric neutralization.
  \item Determine the actual mass of commercial slaked lime ($92.5\%$ purity) that the plant operators must weigh and add to the holding tank.
\end{enumerate}
\end{exercise}

\begin{tabletworkspace}[5.0cm]
\end{tabletworkspace}

\begin{solution}
\textbf{1. Total Acid Normality $N_{\text{total}}$:}
For $\ce{HCl}$, $Z_1 = 1$:
\[
N_1 = Z_1 \cdot C_{M1} = 1 \times 0.0450\ \mathrm{M} = 0.0450\ \mathrm{N}.
\]
For $\ce{H2SO4}$, $Z_2 = 2$:
\[
N_2 = Z_2 \cdot C_{M2} = 2 \times 0.0250\ \mathrm{M} = 0.0500\ \mathrm{N}.
\]
The total acid normality is additive:
\[
N_{\text{total}} = N_1 + N_2 = 0.0450 + 0.0500 = \mathbf{0.0950\ \mathrm{eq\cdot L^{-1}} = 0.0950\ \mathrm{N}}.
\]

\textbf{2. Total Number of Acid Equivalents:}
Total volume is $V = 10.00\ \mathrm{m^3} = 10\,000\ \mathrm{L}$:
\[
\text{Total Equivalents} = N_{\text{total}} \cdot V = 0.0950\ \mathrm{eq\cdot L^{-1}} \times 10\,000\ \mathrm{L} = \mathbf{950.0\ \mathrm{eq}}.
\]

\textbf{3. Balanced Neutralization Equations:}
\begin{align*}
\ce{2 HCl + Ca(OH)2 &-> CaCl2 + 2 H2O} \\
\ce{H2SO4 + Ca(OH)2 &-> CaSO4 v + 2 H2O}
\end{align*}

\textbf{4. Theoretical Mass of Pure $\ce{Ca(OH)2}$:}
Calcium hydroxide provides $2\ \ce{OH-}$ per formula unit ($Z_{\text{base}} = 2$):
\[
E_q(\ce{Ca(OH)2}) = \frac{M}{Z} = \frac{74.09\ \mathrm{g\cdot mol^{-1}}}{2\ \mathrm{eq\cdot mol^{-1}}} = 37.045\ \mathrm{g\cdot eq^{-1}}.
\]
To neutralize $950.0\ \mathrm{eq}$ of acid, exactly $950.0\ \mathrm{eq}$ of base is required:
\[
m_{\text{pure}} = \text{Equivalents} \times E_q = 950.0\ \mathrm{eq} \times 37.045\ \mathrm{g\cdot eq^{-1}} = 35\,192.75\ \mathrm{g} \approx \mathbf{35.19\ \mathrm{kg}}.
\]

\textbf{5. Actual Mass of Commercial Slaked Lime (Purity $92.5\%$):}
\[
m_{\text{commercial}} = \frac{m_{\text{pure}}}{\text{Purity}} = \frac{35.193\ \mathrm{kg}}{0.925} = \mathbf{38.05\ \mathrm{kg}}.
\]
The operators must charge $\mathbf{38.05\ \mathrm{kg}}$ of commercial slaked lime into the holding tank.
\end{solution}
